% MA 214 Discussion 7 -- covers Lectures 9 and 10.
%
% Student handout. Page 1 is the concept review; the question set runs from page 2.
% Compile from inside this folder:  pdflatex Discussion7.tex

\documentclass[11pt]{article}
\usepackage[margin=1in]{geometry}
\usepackage{parskip}
\usepackage{fancyhdr}
\usepackage{array}
\usepackage{booktabs}
\usepackage{enumitem}
\usepackage{amsmath}
\usepackage{tabularx}
\usepackage{listings}

\pagestyle{fancy}
\fancyhf{}
\cfoot{\thepage}
\rhead{Discussion 7}
\lhead{MA 214: Applied Statistics}
\renewcommand{\headrulewidth}{.4mm}

\newcommand{\blankline}[1]{\underline{\hspace{#1}}}
\newcolumntype{L}{>{\raggedright\arraybackslash}X}

\lstset{
  basicstyle=\ttfamily\small,
  columns=fullflexible,
  frame=single,
  framerule=0.4pt,
  breaklines=true,
  showstringspaces=false,
  aboveskip=8pt,
  belowskip=8pt
}

% Section header: bold title over a hairline rule.
\newcommand{\parttitle}[1]{%
  \par\medskip\noindent
  {\large\bfseries #1}\par
  \vspace{-.5em}\noindent\rule{\linewidth}{0.4pt}\par\smallskip}

% Writing space.
\newcommand{\answerspace}[1]{\vspace{#1}}

% Flush-left question list: the label runs in at the margin and the body wraps
% back to the margin, so nothing is pushed far to the right.
\newlist{questions}{enumerate}{1}
\setlist[questions]{label=\textbf{Question \arabic*.}, wide=0pt, itemsep=1.4em, topsep=.8em}

\newlist{parts}{enumerate}{1}
\setlist[parts]{label=\textbf{(\Alph*)}, leftmargin=1.6em, labelsep=.45em,
  itemsep=.7em, topsep=.5em}

\begin{document}

\noindent\textbf{Name:}\blankline{2.3in}\hfill\textbf{BUID:}\blankline{1.6in}

\begin{center}
  {\Large\bfseries Discussion 7: The Same Interval, Three Ways}
\end{center}

\noindent\textbf{Goals:} \textbf{(1)} build one confidence interval by bootstrap and by formula and
compare them, \textbf{(2)} say what does and does not change an interval's width, \textbf{(3)}
recognize when the two bootstrap recipes disagree and why, \textbf{(4)} connect an interval to last
week's test, \textbf{(5)} tell paired data from two samples.

\parttitle{Part 1: Concept Review}

{\small

\noindent Everything in Part 2 can be answered from this page. Keep it in front of you.

\begin{enumerate}[label=\textbf{\arabic*.}, wide=0pt, itemsep=.45em, topsep=.5em]

\item \textbf{Three roads to a sampling distribution.} Every interval and $p$-value in this module
comes from one of these.

\begin{center}
\renewcommand{\arraystretch}{1.15}
\begin{tabular}{lll}
\toprule
\textbf{Road} & \textbf{How you get it} & \textbf{What it needs} \\
\midrule
Simulate & bootstrap: resample $n$ \emph{with replacement}, $B$ times & a sample that is not tiny \\
Formula & the CLT: $\text{SE} = s/\sqrt{n}$ & $n$ large enough for the CLT \\
Exact & compute from the model, e.g.\ binomial & the model to be exactly right \\
\bottomrule
\end{tabular}
\end{center}

\item \textbf{The bootstrap.} Draw a resample of size $n$ \emph{with replacement} from your sample,
compute the statistic, repeat $B$ times. The spread of those $B$ values estimates the sampling
distribution's spread. Resampling \emph{without} replacement just returns your data, and resampling
a size other than $n$ estimates the wrong sampling distribution.

\item \textbf{Two recipes for a bootstrap interval.} With $\hat\theta$ the observed statistic:
\[
\textbf{percentile: } \big[q_{0.025},\, q_{0.975}\big] \text{ of the bootstrap values},
\qquad
\textbf{standard error: } \hat\theta \pm 2\,\text{SE}_{\text{boot}} .
\]
They nearly coincide when the bootstrap distribution is symmetric. When it is skewed, prefer the
percentile interval, because the SE recipe forces symmetry that is not there.

\item \textbf{The formula interval.} For a mean,
\[
\bar{x} \pm t^{\ast}\,\frac{s}{\sqrt{n}},
\qquad \text{df} = n-1 .
\]
Use $t^{\ast}$, not $z^{\ast}$, whenever $\sigma$ is estimated by $s$. The two are close for large
$n$ ($t^{\ast}_{59} = 2.001$ against $z^{\ast} = 1.96$) and far apart for small $n$
($t^{\ast}_{7} = 2.365$, about 21\% wider).

\item \textbf{What changes the width.} Raising $n$ narrows it like $1/\sqrt{n}$, so four times the
data halves the width. Raising the confidence level widens it. Raising $B$ does \textbf{not} narrow
it: more simulations only pin down the interval you already had.

\item \textbf{Interval and test are two views of one thing.} A 95\% interval for a difference that
excludes $0$ corresponds to rejecting $H_0: \text{difference} = 0$ at the 5\% level. The test says
whether; the interval says how big, and in what units.

\item \textbf{Paired or two samples.} If each number in one group is tied to a specific number in
the other, the data are \textbf{paired}: reduce to one difference per pair and analyze those
differences. Treating paired data as two independent samples throws away the pairing and usually
overstates the uncertainty.

\end{enumerate}

}

\newpage

\parttitle{Part 2: Question Set}

\noindent\textbf{Scenario.} For the $n = 60$ library branches, annual visits (in thousands) have
$\bar{x} = 75.202$ and $s = 29.709$. A bootstrap with $B = 2000$ resamples gives a bootstrap
standard error of $3.695$, bootstrap percentiles $q_{0.025} = 67.91$ and $q_{0.975} = 82.55$, and a
bootstrap distribution whose mean ($75.167$) and median ($75.125$) are nearly equal. Useful values:
$t^{\ast}_{59} = 2.001$, \; $z^{\ast} = 1.96$, \; $\sqrt{60} = 7.746$.

\begin{questions}

%%%%%%%%%%%%%%%%%%%%%%%%%%%%%%%%%%%%%%
\item \textbf{Three roads, one interval.}

\begin{parts}
\item Compute the formula standard error $s/\sqrt{n}$, then the 95\% $t$ interval for the mean.
\answerspace{4.5em}

\item Write down the bootstrap percentile interval, and compute the bootstrap standard-error
interval $\bar{x} \pm 2\,\text{SE}_{\text{boot}}$.
\answerspace{4em}

\item You now have three intervals. Fill in their widths and comment on how close they are.

\begin{center}
\renewcommand{\arraystretch}{1.6}
\begin{tabular}{lcc}
\toprule
\textbf{Method} & \textbf{Interval} & \textbf{Width} \\
\midrule
$t$ formula & \blankline{1.7in} & \blankline{.9in} \\
Bootstrap percentile & \blankline{1.7in} & \blankline{.9in} \\
Bootstrap standard error & \blankline{1.7in} & \blankline{.9in} \\
\bottomrule
\end{tabular}
\end{center}

\item The two bootstrap recipes agree to within about $0.1$. Which fact given in the scenario tells
you in advance that they would?
\answerspace{3.5em}

\item The formula SE and the bootstrap SE are $3.835$ and $3.695$. Should you be troubled that they
are not identical? Explain what each one is.
\answerspace{4em}
\end{parts}

%%%%%%%%%%%%%%%%%%%%%%%%%%%%%%%%%%%%%%
\item \textbf{What moves the width.}

\begin{parts}
\item The city funds a survey of $240$ branches instead of $60$, with the same $s$. Compute the new
SE and the new interval width, and state the general rule you just used.
\answerspace{4.5em}

\item A colleague reruns the bootstrap with $B = 100{,}000$ instead of $2000$ and expects a narrower
interval. What actually happens, and what does raising $B$ buy?
\answerspace{4em}

\item Write a one-sentence interpretation of your interval from Question 1 that mentions the
\emph{procedure} and \emph{repeated samples}.
\answerspace{4em}

\item All 60 branches are in one city. Does any of the three methods let you claim an interval for
mean visits at branches nationwide? What would?
\answerspace{4em}
\end{parts}

%%%%%%%%%%%%%%%%%%%%%%%%%%%%%%%%%%%%%%
\item \textbf{When the recipes disagree.} The director asks about the \emph{busiest} branches, so
the statistic is the $90$th percentile of visits. The observed value is $116.87$.
Bootstrapping it gives a bootstrap standard error of $6.681$, with percentiles
$q_{0.025} = 99.43$ and $q_{0.975} = 123.70$.

\begin{parts}
\item Write the percentile interval, then the standard-error interval
$116.87 \pm 2(6.681)$. Compare them.
\answerspace{4.5em}

\item The percentile interval reaches $17.4$ below the estimate but only $6.8$ above it. What does
that asymmetry say about the bootstrap distribution of this statistic?
\answerspace{4em}

\item The two intervals disagree here but agreed in Question 1. Which do you report for the $90$th
percentile, and why?
\answerspace{4em}

\item Which of the three roads in the review is unavailable for this statistic, and what does that
tell you about when the bootstrap earns its keep?
\answerspace{4em}
\end{parts}

%%%%%%%%%%%%%%%%%%%%%%%%%%%%%%%%%%%%%%
\item \textbf{Back to last week's test.} The Saturday-hours comparison from Discussion 6 had
$12$ program branches and $12$ controls, with an observed difference in mean weekly visits of
$0.983$ hundreds. The permutation test gave a two-sided $p = 0.024$. Now estimate the effect:

\begin{center}
\renewcommand{\arraystretch}{1.2}
\begin{tabular}{ll}
\toprule
$t$ interval for the difference & $[\,0.149,\; 1.817\,]$ \\
Bootstrap percentile interval & $[\,0.225,\; 1.734\,]$ \\
\bottomrule
\end{tabular}
\end{center}

\begin{parts}
\item Both intervals exclude $0$, and last week's $p$ was $0.024$. State the general connection
between those two facts.
\answerspace{3.5em}

\item Interpret the $t$ interval for the director, in visits per week rather than hundreds.
\answerspace{4em}

\item The test told the director something the interval does not, and the interval tells her
something the test does not. Say what each contributes.
\answerspace{4em}

\item Now suppose instead that the \emph{same} 12 branches had been measured before and after the
program, with no control group. What would the data be, what would you analyze, and why is the
two-sample interval above no longer the right tool?
\answerspace{4.5em}
\end{parts}

%%%%%%%%%%%%%%%%%%%%%%%%%%%%%%%%%%%%%%
\item \textbf{Find the bugs.} A classmate wants a 95\% bootstrap interval for mean visits. The
vector \texttt{v} holds all $n = 60$ values. The code runs, prints two plausible numbers, and is
wrong in \textbf{three} separate ways.

\begin{lstlisting}[language=R]
boot <- replicate(2000,
  mean(sample(v, size = 30, replace = FALSE)))

quantile(v, c(0.025, 0.975))
\end{lstlisting}

\begin{parts}
\item Bug 1 is \texttt{replace = FALSE}. What does each resample then consist of, and what does
that do to the $2000$ values in \texttt{boot}?
\answerspace{4em}

\item Bug 2 is \texttt{size = 30}. Which sampling distribution is this estimating, and would the
resulting interval be too wide or too narrow?
\answerspace{4em}

\item Bug 3 is the last line. It reports $[29.39,\; 136.61]$, which is about $7.3$ times as wide as
the correct interval. Say precisely what quantity that line has taken the percentiles of, and what
it should have taken them of.
\answerspace{4.5em}

\item Rewrite the block correctly.
\answerspace{4em}

\item The same classmate reports the correct interval $[\,67.5,\; 82.9\,]$ and then writes the two
sentences below. Mark each \textbf{defensible} or \textbf{not}, and rewrite the ones that are not.
Your answer to bug 3 is a useful check on the first one.
\begin{enumerate}[label=\textbf{S\arabic*.}, leftmargin=2.2em, itemsep=.25em, topsep=.35em]
\item ``95\% of branches have annual visits between $67.5$ and $82.9$ thousand.''
\item ``There is a 95\% probability that $\mu$ lies between $67.5$ and $82.9$.''
\end{enumerate}
\answerspace{5em}
\end{parts}

\end{questions}

\end{document}
